%% start of file `template.tex'.
%% Copyright 2006-2013 Xavier Danaux (xdanaux@gmail.com).
%
% This work may be distributed and/or modified under the
% conditions of the LaTeX Project Public License version 1.3c,
% available at http://www.latex-project.org/lppl/.


\documentclass[10pt,a4paper,sans]{moderncv}        % possible options include font size ('10pt', '11pt' and '12pt'), paper size ('a4paper', 'letterpaper', 'a5paper', 'legalpaper', 'executivepaper' and 'landscape') and font family ('sans' and 'roman')

% moderncv themes
\moderncvstyle{classic}                             % style options are 'casual' (default), 'classic', 'oldstyle' and 'banking'
\moderncvcolor{black}                               % color options 'blue' (default), 'orange', 'green', 'red', 'purple', 'grey' and 'black'
%\renewcommand{\familydefault}{\sfdefault}         % to set the default font; use '\sfdefault' for the default sans serif font, '\rmdefault' for the default roman one, or any tex font name
%\nopagenumbers{}                                  % uncomment to suppress automatic page numbering for CVs longer than one page

% character encoding
%\usepackage[utf8]{inputenc}                       % if you are not using xelatex ou lualatex, replace by the encoding you are using
%\usepackage{CJKutf8}                              % if you need to use CJK to typeset your resume in Chinese, Japanese or Korean

% adjust the page margins
\usepackage[scale=0.75]{geometry}
\setlength{\hintscolumnwidth}{3cm}                % if you want to change the width of the column with the dates
%\setlength{\makecvtitlenamewidth}{10cm}           % for the 'classic' style, if you want to force the width allocated to your name and avoid line breaks. be careful though, the length is normally calculated to avoid any overlap with your personal info; use this at your own typographical risks...

% personal data
\name{Yuma}{Mizuno}
\colorlet{firstnamecolor}{lastnamecolor}
%\title{}                               % optional, remove / comment the line if not wanted
\address{School of Mathematical Sciences, University College Cork}{Western Gateway Building, Western Road, Cork, Ireland}% optional, remove / comment the line if not wanted; the "postcode city" and "country" arguments can be omitted or provided empty
%\phone[mobile]{+1~(234)~567~890}                   % optional, remove / comment the line if not wanted; the optional "type" of the phone can be "mobile" (default), "fixed" or "fax"
%\phone[fixed]{+2~(345)~678~901}
%\phone[fax]{+3~(456)~789~012}
\email{mizuno.y.aj@gmail.com}                              % optional, remove / comment the line if not wanted
%\homepage{www.johndoe.com}                         % optional, remove / comment the line if not wanted
%\social[linkedin]{john.doe}                        % optional, remove / comment the line if not wanted
%\social[twitter]{jdoe}                             % optional, remove / comment the line if not wanted
%\social[github]{jdoe}                              % optional, remove / comment the line if not wanted
\extrainfo{\href{https://github.com/yuma-mizuno}{GitHub}\quad\href{https://researchmap.jp/yuma_mizuno}{researchmap}}
%\photo[64pt][0.4pt]{picture}                       % optional, remove / comment the line if not wanted; '64pt' is the height the picture must be resized to, 0.4pt is the thickness of the frame around it (put it to 0pt for no frame) and 'picture' is the name of the picture file
%\quote{Some quote}                                 % optional, remove / comment the line if not wanted

% to show numerical labels in the bibliography (default is to show no labels); only useful if you make citations in your resume
%\makeatletter
%\renewcommand*{\bibliographyitemlabel}{\@biblabel{\arabic{enumiv}}}
%\makeatother
%\renewcommand*{\bibliographyitemlabel}{[\arabic{enumiv}]}% CONSIDER REPLACING THE ABOVE BY THIS

% bibliography with mutiple entries
%\usepackage{multibib}
%\newcites{book,misc}{{Books},{Others}}
\usepackage[backend=biber,style=numeric,sorting=ydnt,defernumbers=true,maxbibnames=99,giveninits=true]{biblatex}
\DeclareFieldFormat*{title}{\mkbibemph{#1}}
\DeclareFieldFormat{journaltitle}{#1}
\DeclareFieldFormat{booktitle}{#1}
\renewbibmacro{in:}{}
\addbibresource{publications.bib}

\defbibenvironment{bibliography}
  {\list{}{
      \setlength{\leftmargin}{0pt}
      \setlength{\itemindent}{0pt}
      \setlength{\itemsep}{\bibitemsep}
      \setlength{\parsep}{\bibparsep}}}
  {\endlist}
  {\item}
\AtEveryBibitem{%
  \ifkeyword{publication}{%
    \clearfield{eprint}%
    \clearfield{eprintclass}%
    \clearfield{eprinttype}%
  }{}
}
%----------------------------------------------------------------------------------
%            content
%----------------------------------------------------------------------------------
\begin{document}
%\begin{CJK*}{UTF8}{gbsn}                          % to typeset your resume in Chinese using CJK
%-----       resume       ---------------------------------------------------------
\makecvtitle

\section{Education}
\cventry{Mar 2021}{Doctor of Science}{Department of Mathematical and Computing Science, Tokyo Institute of Technology}{}{}{Thesis: Difference equations arising from cluster algebras\\ Supervisor: Yuji Terashima and Sakie Suzuki}  % arguments 3 to 6 can be left empty
\cventry{Mar 2018}{Master of Science}{Department of Mathematical and Computing Science, Tokyo Institute of Technology}{}{}{Supervisor: Yuji Terashima}
\cventry{Mar 2016}{Bachelor of Science}{Department of Information Science, Tokyo Institute of Technology}{}{}{}
%\cventry{year--year}{degree or job title}{institution or employer}{city}{grade}{description}

\section{Teaching}
\cvitem{Semester 2, 2025-2026}{Lecturer, University College Cork, Measure Theory and Integration (MA3064)}
\cvitem{Fall 2021}{Part-time lecturer, Tokyo University of Agriculture and Technology, Calculus II and Exercises}

\section{Research Positions and Fellowships}
\cvitem{2025--present}{Postdoctoral Researcher, University College Cork}
\cvitem{2024--2025}{Postdoctoral Researcher, University College Dublin}
\cvitem{2021--2024}{JSPS Research Fellow (PD), Japan Society for the Promotion of Science}
\cvitem{2018--2021}{JSPS Research Fellow (DC1), Japan Society for the Promotion of Science}


%\section{Languages}
%\cvitemwithcomment{Language 1}{Skill level}{Comment}
%\cvitemwithcomment{Language 2}{Skill level}{Comment}
%\cvitemwithcomment{Language 3}{Skill level}{Comment}
%
%\section{Computer skills}
%\cvdoubleitem{category 1}{XXX, YYY, ZZZ}{category 4}{XXX, YYY, ZZZ}
%\cvdoubleitem{category 2}{XXX, YYY, ZZZ}{category 5}{XXX, YYY, ZZZ}
%\cvdoubleitem{category 3}{XXX, YYY, ZZZ}{category 6}{XXX, YYY, ZZZ}

%\section{Interests}
%\cvitem{hobby 1}{Description}
%\cvitem{hobby 2}{Description}
%\cvitem{hobby 3}{Description}

%\section{Extra 1}
%\cvlistitem{Item 1}
%\cvlistitem{Item 2}
%\cvlistitem{Item 3. This item is particularly long and therefore normally spans over several lines. Did you notice the indentation when the line wraps?}
%
%\section{Extra 2}
%\cvlistdoubleitem{Item 1}{Item 4}
%\cvlistdoubleitem{Item 2}{Item 5\cite{book1}}
%\cvlistdoubleitem{Item 3}{Item 6. Like item 3 in the single column list before, this item is particularly long to wrap over several lines.}
%
%\section{References}
%\begin{cvcolumns}
%  \cvcolumn{Category 1}{\begin{itemize}\item Person 1\item Person 2\item Person 3\end{itemize}}
%  \cvcolumn{Category 2}{Amongst others:\begin{itemize}\item Person 1, and\item Person 2\end{itemize}(more upon request)}
%  \cvcolumn[0.5]{All the rest \& some more}{\textit{That} person, and \textbf{those} also (all available upon request).}
%\end{cvcolumns}
\nocite{*}
\begin{refcontext}[sorting=none]
\printbibliography[title={Preprints},keyword={preprint}]
\end{refcontext}
\printbibliography[title={Publications},keyword={publication}]

\section{Invited Talks}

\cvitem{13 Jul 2026}{\textbf{Nahm sums from cluster transformations}, \textit{Workshop on Higher Teichmüller theory, knot theory and cluster algebras}, Cork, Ireland.}
\cvitem{18 May 2026}{\textbf{Adjoint Reidemeister torsion for $G$-local systems on 3-manifolds with torus boundary}, \textit{Paris algebra seminar}, Paris, France.}
\cvitem{15 Apr 2026}{\textbf{Adjoint Reidemeister torsion for $G$-local systems on 3-manifolds with torus boundary}, \textit{Research Seminar on Quantum Topology and Categorification (QTCat)}, Hamburg, Germany.}
\cvitem{17 Feb 2026}{\textbf{Formalizing mathematics with Lean theorem prover}, \textit{Boole Colloquium in Mathematics}, Cork, Ireland.}
\cvitem{9 Dec 2025}{\textbf{Adjoint Reidemeister torsion for $G$-local systems on 3-manifolds}, \textit{Geometry and Mathematical Physics seminar}, Birmingham, UK.}
\cvitem{13 Nov 2025}{\textbf{Formalizing mathematics with theorem provers}, \textit{The 28th Workshop on Information-Based Induction Sciences (IBIS2025)}, Naha, Japan.}
\cvitem{31 Oct 2025}{\textbf{Introduction to Lean theorem prover}, \textit{iTHEMS Math \& Computer Seminar}, Wako, Japan.}
\cvitem{18 Jun 2025}{\textbf{Lean theorem prover and monoidal category theory}, \textit{BIMSA-Tsinghua Quantum Symmetry Seminar}, Online.}
\cvitem{19 Feb 2025}{\textbf{Geometry of $q$-Painlevé equations and cluster structures}, \textit{Asymptotic Expansion of $\tau$-functions and Related Topics}, Kyoto, Japan.}
\cvitem{12 Mar 2024}{\textbf{Cluster algebras and $q$-series}, \textit{$q$-Series and Related Topics}, Osaka, Japan.}
\cvitem{11 Mar 2024}{\textbf{Cluster structures on $q$-Painlevé systems via toric geometry}, \textit{Advances in Cluster Algebras 2024}, Nagoya, Japan.}
\cvitem{16 Dec 2023}{\textbf{Geometry of $q$-Painlevé systems and cluster structures}, \textit{Comprehensive Studies of Differential Equations}, Tokyo, Japan.}
\cvitem{21 Jun 2023}{\textbf{Classification of periodic $Y$-systems of rank 2}, \textit{Cluster algebras and related topics in East Asia}, Seoul, South Korea.}
\cvitem{30 May 2023}{\textbf{Modularity of Nahm sums and periodicity phenomena in cluster algebras}, \textit{South Osaka Algebra Seminar}, Osaka, Japan.}
\cvitem{26 Oct 2022}{\textbf{$q$-Painlevé systems, toric surfaces, and cluster Poisson varieties}, \textit{Painlevé Equations: From Classical to Modern Analysis}, Strasbourg, France.}
\cvitem{14 May 2022}{\textbf{Nahm's problem and cluster algebras}, \textit{AMS Spring Western Virtual Sectional Meeting}, Online.}
\cvitem{18 Nov 2021}{\textbf{Nahm's problem and cluster algebras}, \textit{UCD Algebra and Number Theory seminar}, Online.}
\cvitem{28 Sep 2021}{\textbf{$q$-Painlevé systems on cluster Poisson varieties}, \textit{Infinite Analysis 21 Workshop Around Cluster Algebras}, Online.}
\cvitem{17 Jul 2020}{\textbf{$T$-systems and $Y$-systems in cluster algebras}, \textit{South Osaka Algebra Seminar}, Online.}
\cvitem{18 Jun 2020}{\textbf{Difference equations arising from cluster algebras}, \textit{Representation Theory Seminar (RIMS)}, Online.}
\cvitem{12 Dec 2019}{\textbf{The Cartan--Killing classification and its extensions in cluster algebras}, \textit{Workshop on Combinatorics and Related Topics in Shinshu, Winter 2019}, Matsumoto, Japan.}
\cvitem{13 Feb 2019}{\textbf{Exponents associated with $Y$-systems}, \textit{Tokyo Institute of Technology Topology Seminar}, Tokyo, Japan.}
\cvitem{22 Jan 2019}{\textbf{Combinatorics of ideal triangulations and quiver mutations}, \textit{The 14th East Asian Conference on Geometric Topology}, Beijing, China.}
\cvitem{6 Dec 2017}{\textbf{Ideal triangulations of 3-manifolds and mutation loops}, \textit{Infinite Analysis 17 -Algebraic and Combinatorial Aspects in Integrable Systems-}, Osaka, Japan.}

\section{Other Talks and Presentations}

\cvitem{17 Jan 2025}{\textbf{Metaprogramming on monoidal categories}, \textit{Lean Together 2025}, Online.}
\cvitem{30 Oct 2023}{\textbf{Computations with morphisms in monoidal categories in Lean}, \textit{The 19th Theorem Proving and Provers meeting}, Tokyo, Japan.}
\cvitem{22 Oct 2022}{\textbf{Computations on monoidal categories in the proof assistant Lean}, \textit{Topology and Computers 2022}, Higashihiroshima, Japan.}
\cvitem{14 Sep 2022}{\textbf{Mutations of blowups of toric surfaces and $q$-Painlevé systems}, \textit{Mathematical Society of Japan Autumn Meeting 2022}, Sapporo, Japan.}
\cvitem{11 Aug 2022}{\textbf{Mutations of blowups of toric surfaces and $q$-Painlevé systems}, \textit{Summer Seminar on Functional Equations 2022}, Mori, Japan.}
\cvitem{31 Mar 2022}{\textbf{Mutations of blowups of toric surfaces and $q$-Painlevé system}, \textit{Mathematical Society of Japan Annual Meeting 2022 (abstract only; no oral presentation)}, Online.}
\cvitem{25 Jun 2021}{\textbf{Cluster algebras and spaces of initial conditions for $q$-Painlevé systems}, \textit{ALTReT 2021}, Online.}
\cvitem{4 Mar 2021}{\textbf{$q$-Painlevé systems on cluster varieties}, \textit{The 17th Mathematics Conference for Young Researchers}, Sapporo, Japan.}
\cvitem{17 Jun 2019}{\textbf{Exponents associated with periodic $Y$-systems}, \textit{Cluster Algebras 2019 (CA19) (poster)}, Kyoto, Japan.}
\cvitem{19 Mar 2019}{\textbf{Exponents associated with $Y$-systems of type $X_r$ at level $l$}, \textit{Mathematical Society of Japan Annual Meeting 2019}, Tokyo, Japan.}
\cvitem{6 Mar 2019}{\textbf{Determinant identities associated with mutation sequences}, \textit{The 15th Mathematics Conference for Young Researchers}, Sapporo, Japan.}
\cvitem{24 Sep 2018}{\textbf{Jacobian matrices of $Y$-seed mutations and mutation networks}, \textit{Mathematical Society of Japan Autumn Meeting 2018}, Okayama, Japan.}
\cvitem{13 Feb 2018}{\textbf{Jacobian matrices of cluster transformations and representations of affine Lie algebras}, \textit{The 1st Seminar for New Researchers in Mathematics}, Kyoto, Japan.}
\cvitem{26 Mar 2017}{\textbf{Quiver mutations and $q$-binomial identities}, \textit{Mathematical Society of Japan Annual Meeting 2017}, Hachioji, Japan.}
\end{document}


%% end of file `template.tex'.
